\begin{problem}
    El crecimiento de una población de bacterias se modela utilizando la siguiente expresión:
    \[
    F = P(1 + r)^t,
    \]
    tal que $F$ es la cantidad final de bacterias transcurrido un tiempo $t$ en horas, $P$ es la cantidad inicial de bacterias y $r$ es la tasa de crecimiento por hora.

    Para calcular la cantidad de bacterias que habrá transcurridas \num{2} horas, en un cultivo en que inicialmente hay \num{100} bacterias y tiene una tasa de crecimiento de $\dfrac{1}{2}$ por hora, se realiza el siguiente procedimiento cometiéndose un error.

    \begin{itemize}
        \item Paso 1: se reemplazan los valores de $P$, $t$ y $r$, obteniéndose $F = 100\left(1 + \dfrac{1}{2}\right)^2$.
        \item Paso 2: se calcula la suma dentro del paréntesis, obteniéndose $F = 100\left(\dfrac{3}{2}\right)^2$.
        \item Paso 3: se resuelve la potencia, obteniéndose $F = 100 \cdot \dfrac{6}{4}$.
        \item Paso 4: se resuelve el producto, obteniéndose $F = 150$.
    \end{itemize}

    ¿En cuál de los pasos se cometió el error?

    \begin{enumerate}[label=\Alph*.]
        \item En el Paso 1
        \item En el Paso 2
        \item En el Paso 3
        \item En el Paso 4
    \end{enumerate}
\end{problem}
